\subsection{Stärken}
\label{subsec:staerken}

Die architektonische Hauptstärke ist die Verbindung einer gemeinsamen
Baseline mit expliziten Verantwortungsgrenzen. Frontend, Backend,
Desktop-Schicht, Persistenz und Tooling bleiben getrennte Bausteine; zugleich
wird das Frontend profilübergreifend verwendet. Damit adressiert die
Architektur die Anforderungen A-02 und A-05, ohne jede Variante mit allen
Laufzeitkomponenten zu belasten (Kapitel~\ref{sec:architektur}).

Bei der Variabilität bilden fünf Profile wiederkehrende Plattformzuschnitte
ab. Abhängigkeiten verhindern insbesondere PostgreSQL ohne Backend, während
die Capability aus kompatiblen Profilen herausgelöst bleibt. Die dokumentierte
Abnahme der fünf Basisprofile, der drei zulässigen PostgreSQL-Kombinationen und
der beiden unzulässigen Kombinationen stützt diese Stärke für den aktuellen,
begrenzten Feature-Katalog \parencite{kleiveist2026acceptance}. Sie belegt
nicht, dass beliebig viele künftige Capabilities ebenso beherrschbar wären.

Prozessual schafft \texttt{control.py} einen gemeinsamen Einstieg, ohne npm,
Pytest, Cargo, Tauri, Docker oder Alembic zu verbergen. Profilabhängiges
Überspringen nicht vorhandener Schichten verbindet eine einheitliche
Befehlsoberfläche mit den tatsächlichen Projektgrenzen
(Kapitel~\ref{sec:tooling-entwicklungsworkflow}). Für Qualität und
Nachvollziehbarkeit sprechen das versionierte Profilmanifest, der
Konfigurationsvertrag, Tests und Buildpfade sowie die dokumentierte
byte-identische Generierung zweier \texttt{desktop-cloud + postgres}-Projekte
bei gleichen Eingaben \parencite{kleiveist2026acceptance}. Der Nachweis zur
Reproduzierbarkeit ist damit konkret, aber auf eine untersuchte Kombination
und Umgebung begrenzt.

Tabelle~\ref{tab:strengths-limitations} verdichtet diese Stärken gemeinsam mit
ihren jeweiligen Geltungsgrenzen. Sie trennt beobachtete Mechanismen und
dokumentierte Prüfergebnisse von der daraus abgeleiteten Bedeutung.

\begin{table}[H]
  \centering
  \small
  \renewcommand{\arraystretch}{1.2}
  \caption{Struktur zur Gegenüberstellung von Stärken und Grenzen}
  \label{tab:strengths-limitations}
  \begin{tabularx}{\textwidth}{@{}p{0.18\textwidth}p{0.27\textwidth}p{0.27\textwidth}>{\raggedright\arraybackslash}X@{}}
    \toprule
    \textbf{Aspekt} & \textbf{Stärke / Evidenz} & \textbf{Grenze / Evidenz} & \textbf{Auswirkung} \\
    \midrule
    % Später ausschließlich hergeleitete Stärken und Grenzen eintragen.
    % Beobachtung, Bewertung und Auswirkung nachvollziehbar trennen.
    \bottomrule
  \end{tabularx}
\end{table}

